\section{The self-model is a hardware primitive}
\label{sec:selfmodel}

\subsection{Learning is the third term, not a second loop}

A conventional learning system has two loops: a fast inference loop over a
frozen network, and a slow training loop (back-propagation) that edits the
network from outside. The two use different data paths, different memory, and
often different hardware. \TALOS{} has one loop. The regeneration term
$\Regen$ of the single instruction (Eq.~\ref{eq:lomega}) drives the state
toward its canonical nonlinear anchored self-model
\begin{equation*}
\varphi_J(\Gam) \;=\; \bigl(1 - R(\Gam)\bigr)\,\mathcal{P}_{\mathrm{Fano}}(\Gam)
\;+\; R(\Gam)\,\ket{\psi}\bra{\psi},
\qquad R(\Gam) \;=\; \frac{1}{7\,P(\Gam)},
\end{equation*}
where $\mathcal{P}_{\mathrm{Fano}}$ is the Fano pinching channel
(\S\ref{sec:fabric}) and $\ket{\psi}\bra{\psi}$ is the collineation-invariant
anchor (explicitly distinct from linear relaxation toward a static target
$\rho_{\mathrm{target}}$). Thus \emph{adaptation is a factor of the tick, not a
procedure above it}: there is no separate optimiser to schedule, no gradient
buffer to store, and no train/infer divide.

The self-reference is mathematically well-posed: because $R(\Gam)\in[1/7,1]$
for every $\Gam\in\Dens(\mathbb{C}^7)$, $\varphi_J$ maps the compact convex
state manifold continuously into itself, so fixed-point existence follows by
Brouwer's theorem (with exact unique fixed point
$\rhostar = \tfrac12\tfrac{\mathbb{1}}{7} + \tfrac12\ket{\psi}\bra{\psi}$ for
the uniform anchor; tower convergence requires dissipative backbone dominance
rather than unconditional global Banach contraction, T-191, verified in
\texttt{website/scripts/check\_core\_numbers.py} and
\texttt{architecture/talos\_v0\_emulator.py}). One honest distinction, forced
by the power-port theorem (\S\ref{subsec:powerport}): retention against
dissipation is \emph{powered} by the non-unital negentropy supply (through
$\ket{\psi}\bra{\psi}$ / the power port) and \emph{directed} by $\Regen$; the
identity below concerns where the imported order goes, not where the free
energy comes from.

\begin{zerobox}[From zero --- the mirror game]
The toy carries a small portrait of itself, painted by itself. Every shake of
the globe does one extra thing: it compares toy to portrait and nudges the
toy a little toward the better version in the picture. That is the whole of
``learning'' here --- no teacher, no lesson-time, no separate school: the
nudge is the third stride of every patrol (Figure~\ref{fig:tickwheel}). And
because the portrait is painted from the inside, keeping the toy \emph{alive}
and making the toy \emph{better} turn out to be the same nudge --- the
identity proved next.
\end{zerobox}

\subsection{A second new identity: refresh is learning}

Dynamic memory must be refreshed --- leaked charge is periodically rewritten
to its stored value, at a continuous energy cost that computes nothing. In
\TALOS{} the analogue of refresh is the regeneration $\Regen$ holding $\Gam$
near $\rhostar$ against the dissipator's leak. But $\rhostar$ is not a stored
value: it is the state's \emph{self-model}. Hence:

\begin{principle}[Refresh--learn identity]\label{prin:refresh}
Retention and adaptation are the same operation. The process that keeps a
holon's state alive (regeneration toward $\varphi_J(\Gam)$) is the process that
adapts (moves toward the self-model). There is no cost that maintains without
also adapting, and no adaptation that is not also maintenance. DRAM's
refresh energy --- spent to compute nothing --- has no counterpart: every
maintenance tick is an adaptation tick.
\end{principle}

\noindent Together with the compute--diagnose identity
(Principle~\ref{prin:selfdiag}), this collapses a \emph{second} pair of
traditionally-separate subsystems: memory refresh and model training, which
in a layered machine are unrelated, are one term here.

\subsection{The reflection threshold as a state-space window bound}

In UHM, the reflection measure is $R(\Gam) = 1/(7\,P(\Gam))$. The L2 reflection
gate $R(\Gam) \ge \Rth = 1/3$ is therefore \textbf{algebraically equivalent}
(\status{T}) to the upper purity boundary of the viability window:
\[
R(\Gam) \;=\; \frac{1}{7\,P(\Gam)} \;\ge\; \frac{1}{3}
\;\Longleftrightarrow\;
P(\Gam) \;\le\; \frac{3}{7}.
\]
Notice what this theorem does and does \emph{not} say: $\Rth=1/3$ is a
\emph{state-space bound} ensuring the holon remains inside the viability window
$P\in(2/7,3/7]$ (avoiding over-purification above $3/7$ as well as decoherence
below $2/7$), \emph{not} a scheduler rule prescribing that $1/3$ of clock
cycles be spent executing $\Regen$. Together with the exact integration gate
$\Phi=(P-P_{\mathrm{diag}})/P_{\mathrm{diag}}\ge 1$ and differentiation gate
$D_{\mathrm{diff}}=1+6\,\CohE\ge 2$, these thresholds are monitored directly as
per-tick state observables (\S\ref{sec:abi}). Any numerical step-size damping
chosen near the viability boundary is an implementation/control choice
(\status{D}) rather than a derived theorem, and any correspondence between
$\Phi$ crossing $1$ and delayed generalisation (``grokking'') in neural
networks remains an open empirical hypothesis (\status{H}).
